\section{Thực nghiệm và Demo minh hoạ}
\label{sec:thuc-nghiem}

% ============================================================
% MỤC TIÊU CHƯƠNG:
% Minh chứng thực nghiệm cho lý thuyết ở Chương 3 & 4:
% equivariant model cần ít data hơn, ít parameters hơn,
% nhưng cho kết quả tốt hơn.
% ============================================================

\subsection{Bài toán và Tập dữ liệu}
\label{subsec:dataset}
% Trình bày:
%   - Giới thiệu bài toán được chọn
%     (gợi ý: QM9 molecular property prediction, N-body simulation,
%     hoặc Rotated MNIST nếu muốn đơn giản hơn)
%   - Mô tả dataset: kích thước, đặc điểm, tính đối xứng có trong data
%   - Lý do chọn: thể hiện rõ lợi ích của equivariance

\subsection{Mô tả mô hình cài đặt}
\label{subsec:mo-hinh}
% Trình bày 2 mô hình so sánh:
%
%   Mô hình Baseline (không có tính đối xứng):
%   - Kiến trúc: MLP thông thường / Standard GNN
%   - Không encode bất kỳ symmetry nào
%
%   Mô hình Đề xuất (equivariant):
%   - Một trong các lựa chọn:
%     * DeepSets (permutation invariant, đơn giản nhất)
%     * E(n)-Equivariant GNN / EGNN (Satorras et al ICML 2021)
%     * Frame Averaging (Puny et al 2022) áp dụng lên MLP bình thường
%   - Nêu rõ điểm mở rộng so với code gốc
%
%   Lưu ý: Liên kết rõ với Chương 3
%   - Nếu dùng DeepSet → dẫn về mục 3.1.4 (parameter sharing)
%   - Nếu dùng Frame Averaging → dẫn về mục 3.1.3 (Paper 3 Phúc)
%   - Nếu dùng EGNN → dẫn về mục 3.2.5 (NequIP, Paper 2 Thịnh)

\subsection{Môi trường thực nghiệm}
\label{subsec:moi-truong}
% Trình bày:
%   - Ngôn ngữ: Python + PyTorch / PyTorch Geometric
%   - Phần cứng: GPU / Cloud (ghi rõ)
%   - Cấu trúc thư mục, requirements.txt, README.md

\subsection{Kết quả và Phân tích}
\label{subsec:ket-qua}
% Trình bày:
%   - Bảng/đồ thị so sánh: Loss, Accuracy, MAE
%   - Chứng minh data efficiency:
%     Equivariant model đạt cùng accuracy với ÍT data hơn
%   - So sánh số lượng parameters
%   - Phân tích kết quả: tại sao equivariant model tốt hơn?
%     (liên kết về lý thuyết Chương 4.1)
%
%   Gợi ý bonus (+2 điểm demo ấn tượng):
%   Vẽ thêm: OOD generalization curve
%   (test trên rotated/permuted data mà chưa thấy trong train)

\clearpage